% Plantilla editorial IMTC
\documentclass[spanish,letterpaper,oneside]{article}

\def\documenttitle {Procesos De Manufactura}
\def\documentsubtitle {Reporte integral - procesos-de-manufactura-imtc}
\def\documentsubject {Ingenieria en Mecatronica}
\def\documentauthor {Martin Jonathan de la Cruz}
\def\coursename {Procesos De Manufactura}
\def\coursecode {IMTC}
\def\universityname {Universidad Autónoma de Nuevo León}
\def\universityfaculty {Facultad de Ingeniería Mecánica y Eléctrica}
\def\universitydepartment {\coursename}
\def\universitydepartmentimage {UANL/assets-uanl/logo-uanl.png}
\def\universitydepartmentimagecfg {width=3.2cm}
\def\universitylocation {Ciudad Universitaria, San Nicolás de los Garza, Nuevo León, México}
\def\authortable {
  \begin{tabular}{ll}
    \textbf{Alumno:} & Martin Jonathan de la Cruz \\
    Matricula: & UANL-IMTC \\
    Programa: & Ingenieria en Mecatronica \\
    Asignatura: & \coursename \\
    Codigo: & \coursecode \\
    Figura docente: & Nombre por definir \\
    & \\
    \multicolumn{2}{l}{Fecha de realizacion: \today} \\
    \multicolumn{2}{l}{\universitylocation}
  \end{tabular}
}

\input{template}

\begin{document}
\templatePortrait
\templatePagecfg
\onehalfspacing

\begin{abstractd}
\textbf{Objetivo:} Documentar el avance academico y tecnico de la materia procesos de manufactura.
\end{abstractd}

\templateIndex
\templateFinalcfg

\section{Introduccion}
Describir el contexto de la actividad y su impacto en la formacion mecatronica.

\section{Desarrollo}
Agregar analisis, resultados, figuras y tablas segun la actividad.

\section{Conclusiones}
Sintetizar hallazgos y siguientes acciones de aprendizaje.

\bibliographystyle{natnumurl}
\bibliography{procesos-de-manufactura-imtc}

\end{document}
